\documentclass[11pt,a4paper]{article}
\usepackage[T1]{fontenc}\usepackage[utf8]{inputenc}
\usepackage{lmodern,microtype,amsmath,amssymb,amsthm,mathtools}
\usepackage[margin=25mm]{geometry}
\usepackage{booktabs,array}\usepackage[hyphens]{url}\usepackage[hidelinks]{hyperref}
\setlength{\emergencystretch}{3em}\allowdisplaybreaks
\newtheorem{theorem}{Theorem}[section]
\newtheorem{lemma}[theorem]{Lemma}
\newtheorem{proposition}[theorem]{Proposition}
\newcommand{\ph}{\Phi}
\newcommand{\ord}{\operatorname{ord}}
\hypersetup{pdftitle={A block reduction for Karmakar--Jana Conjecture 1: proof candidate and method-duplication audit},pdfauthor={Reiwa}}
\title{A block reduction for Karmakar--Jana Conjecture 1\\[5pt]\large A global proof candidate with a failed new-method gate}
\author{Reiwa --- Tournament round 7}
\date{16 November 2026}
\begin{document}\maketitle
\begin{abstract}
We examine Conjecture 1 in Karmakar and Jana, arXiv:2610.11778v1, which
strengthens a two-parameter shifted $q$-supercongruence from
$\Phi_m(q)^3$ to $[m]\Phi_m(q)^2$ for every odd $m>1$ and every permitted
shift. A signed-residue block reduction and a termwise valuation bound
provide a proposed proof over $\mathbb Q(c,d)(q)$. Independent verification
is pending. Crucially, the authors' later paper arXiv:2610.15050 already
uses the same root-of-unity block strategy for a different problem.
We therefore do not claim a substantively new general method or satisfaction
of the tournament's G5 gate. The initially suggested Ni Conjecture 5.1 is
excluded because the original author's version 2 already proves it.
\end{abstract}

\section{Exact source and original assertion}
Let $m>1$ be odd, $h=(m-1)/2$, and $0\leq s\leq h$.
Write
\[
 (a;Q)_k=\prod_{j=0}^{k-1}(1-aQ^j),\qquad
 [j]=\frac{1-q^j}{1-q},\qquad (a,b;Q)_k=(a;Q)_k(b;Q)_k.
\]
We work with independent indeterminates $c,d$ over $K=\mathbb Q(c,d)$.
The original left side is
\begin{equation}\label{eq:L}
 L_{m,s}(q)=\sum_{k=s}^{h}[4k+1]
 \frac{(q;q^2)_{k-s}(q;q^2)_{k+s}(q;q^2)_k^2(cq,dq;q^2)_k}
 {(q^2;q^2)_{k-s}(q^2;q^2)_{k+s}(q^2;q^2)_k^2(q^2/c,q^2/d;q^2)_k}
 (cd)^{-k}q^k.
\end{equation}
For a compact but exact statement of the right side define
\begin{gather}
 A_k=(q;q^2)_k, \qquad C_k=(q^2;q^2)_k, \qquad R_k=A_k/C_k,\label{eq:notation}\\
 P_k=\frac{(cq,dq;q^2)_k}{(q^2/c,q^2/d;q^2)_k}(cd)^{-k}q^k.\notag
\end{gather}
where all these functions depend on $q$ (and $P_k$ also on $c,d$). Set
\begin{align}
 B_{m,s}(q)&=[m+2s]R_s^2P_s q^{-h+s}\frac{R_{h+s}}{R_{h-s}},\label{eq:B}\\[3pt]
 U_{s,j}(q)&=
 \frac{(q,q/(cd);q^2)_j(q^{2s+1};q^2)_j^2}
 {(q^2;q^2)_j^2(q^{2s+2}/c,q^{2s+2}/d;q^2)_j}q^{2j},\label{eq:U}\\[3pt]
 V_{m,s}(q)&=B_{m,s}(q)\sum_{j=0}^{h-s}U_{s,j}(q).\label{eq:V}
\end{align}
The factor $P_s q^{-h+s}$ contains precisely
$(cd)^{-s}q^{-h+2s}(cq,dq;q^2)_s/(q^2/c,q^2/d;q^2)_s$.
Thus \eqref{eq:V} is exactly the right side of equation (1.14), page 3,
of Karmakar--Jana~\cite{KJ}, without changing any index or parameter.
Their Theorem 1 proves
\begin{equation}\label{eq:known}
 L_{m,s}\equiv V_{m,s}\pmod{\Phi_m(q)^3}.
\end{equation}
Their Conjecture 1, Section 4, page 11, asks for the stronger modulus
\begin{equation}\label{eq:target}
 L_{m,s}\equiv V_{m,s}\pmod{[m]\Phi_m(q)^2}.
\end{equation}
The two original pages were checked visually. The latest version shown
by the official arXiv record during this review is v1, 8 October 2026.
No assertion concerning the separate Conjecture 2 is made here.

Congruences are in the reduced-denominator sense over $K[q]$. If a
rational function is regular at every root of a modulus and its reduced
numerator is divisible by that modulus, the congruence holds. Negative
powers of $q$ are harmless because every cyclotomic factor used here has
nonzero roots. Formal parameter congruences do \emph{not} automatically
survive a specialization that makes their reduced denominator vanish at
a cyclotomic root. We do not assert such an unrestricted specialization
principle.

\begin{theorem}[Candidate for the formal-parameter assertion]\label{thm:main}
For every odd integer $m>1$ and every integer $0\leq s\leq(m-1)/2$,
\eqref{eq:target} holds over $K=\mathbb Q(c,d)$.
\end{theorem}
The argument below uses the already proved \eqref{eq:known}, not an
unproved auxiliary conjecture. It supplies all missing proper-divisor
cyclotomic factors. It remains a submitted proof candidate pending
independent checking.

\section{Local orders and the right side}
Fix an odd divisor $\ell>1$ of $m$, put $t=(\ell-1)/2$, and let
$\nu_\ell$ denote the order of a rational function at $\Phi_\ell(q)$.
For a nonnegative integer $k=a\ell+r$, $0\leq r<\ell$, define
\[
 \delta(k)=\begin{cases}0,&0\leq r\leq t,\\1,&t<r<\ell.\end{cases}
\]
Counting the factors with exponent divisible by $\ell$ gives
\begin{equation}\label{eq:valuations}
 \nu_\ell(A_k)=\left\lfloor\frac{k+t}{\ell}\right\rfloor,
 \qquad \nu_\ell(C_k)=\left\lfloor\frac{k}{\ell}\right\rfloor,
 \qquad \nu_\ell(R_k)=\delta(k).
\end{equation}
Each zero is simple, and $\gcd(2,\ell)=1$. All parameter factors such as
$1-cq^e$, $1-q^e/c$, and $1-q^e/(cd)$ are units at $\Phi_\ell$ over $K$.
In particular $\nu_\ell(P_k)=0$. The separate factors $R_k$ are all regular.

\begin{lemma}\label{lem:right}
For every $\ell>1$ dividing $m$,
\[
 \nu_\ell(B_{m,s})=1,\qquad \nu_\ell(U_{s,j})\geq0\quad(j\geq0).
\]
Consequently every summand $B_{m,s}U_{s,j}$, and hence $V_{m,s}$, has
order at least one at $\Phi_\ell$.
\end{lemma}
\begin{proof}
Write $s=b_0\ell+u$, $0\leq u<\ell$. Since $m/\ell$ is odd,
$h=A\ell+t$ with $A=(m/\ell-1)/2$. The order of the prefactor is
\[
 \nu_\ell(B_{m,s})=
 \mathbf1_{u=0}+2\delta(s)+\delta(h+s)-\delta(h-s).
\]
Here $\nu_\ell([m+2s])=\mathbf1_{u=0}$, because $\ell$ is odd and
$m+2s>0$. The three residue cases give
\begin{center}
\begin{tabular}{c|cccc|c}
$u$ & $\mathbf1_{u=0}$ & $\delta(s)$ & $\delta(h+s)$ & $\delta(h-s)$ & total\\\hline
$0$ &1&0&0&0&1\\
$1\leq u\leq t$ &0&0&1&0&1\\
$t<u<\ell$ &0&1&0&1&1
\end{tabular}
\end{center}
This also accounts for the possible pole of $1/R_{h-s}$; it is cancelled
by the other factors before evaluating the prefactor.

For $j=a\ell+r$, $0\leq r<\ell$, write
\[
 e=\left\lfloor\frac{u+t+r}{\ell}\right\rfloor
 -\left\lfloor\frac{u+t}{\ell}\right\rfloor\in\{0,1\}.
\]
The shifted factorial $(q^{2s+1};q^2)_j=A_{s+j}/A_s$ has order $a+e$.
The only nonunit factors in $U_{s,j}$ are therefore counted by
\begin{equation}\label{eq:Uorder}
 \nu_\ell(U_{s,j})=(a+\delta(j))+2(a+e)-2a
 =a+\delta(j)+2e\geq0.
\end{equation}
This proves the claim term by term, without relying on cancellation in
the inner sum.
\end{proof}

In particular, for the smaller problem with $m=\ell$ and any
$0\leq v\leq t$, the known theorem \eqref{eq:known} and
Lemma~\ref{lem:right} imply
\begin{equation}\label{eq:smallzero}
 L_{\ell,v}(\zeta)=0
\end{equation}
for every primitive $\ell$th root $\zeta$.
Its left side is regular since all $R_k$ have nonnegative order.
This use of the known theorem is not circular: only its original
$\Phi_\ell^3$ modulus is used, not the conjectural $[m]$ strengthening.

\section{Exact block limits of the left-side factors}
Let $\zeta$ be an arbitrary primitive $\ell$th root, including when
$\ell$ is composite, and let
\[
 \beta_a=\frac{(1/2)_a}{a!}=\prod_{j=0}^{a-1}\frac{2j+1}{2j+2}.
\]
All root evaluations below mean regular limits of rational functions.

\begin{lemma}\label{lem:blocks}
For $k=a\ell+r$, $a\geq0$ and $0\leq r<\ell$,
\begin{equation}\label{eq:Rblock}
 R_k(\zeta)=
 \begin{cases}\beta_a R_r(\zeta),&0\leq r\leq t,\\
 0,&t<r<\ell.
 \end{cases}
\end{equation}
In addition,
\begin{equation}\label{eq:Pblock}
 P_{a\ell+r}(\zeta)=P_r(\zeta),\qquad
 [4(a\ell+r)+1]_\zeta=[4r+1]_\zeta.
\end{equation}
\end{lemma}
\begin{proof}
In the $j$th complete block of the ratio $R_k$, the unique zero numerator
factor has exponent $(2j+1)\ell$, and the unique zero denominator factor
has exponent $(2j+2)\ell$. Their quotient tends to
$(2j+1)/(2j+2)$. The nonzero numerator and denominator residues each run
through all nonzero classes modulo $\ell$, so their products cancel.
The incomplete block has no zero when $r\leq t$, and one uncancelled
numerator zero when $r>t$. This proves \eqref{eq:Rblock} and regularity.

For a single parameter $z=c$ or $d$, a complete block contributes
\[
 \prod_{j=0}^{\ell-1}\frac{1-z\zeta^{2j+1}}{1-\zeta^{2j+2}/z}
 z^{-\ell}
 =\frac{1-z^\ell}{1-z^{-\ell}}z^{-\ell}=-1.
\]
This is a rational identity in $z$; $c,d$ remain indeterminates. The two
parameters yield a product $(-1)^2=1$, while $q^\ell$ evaluates to one.
Thus $P_k$ is periodic as in \eqref{eq:Pblock}. The final identity follows
directly from $\zeta^\ell=1$.
\end{proof}

\section{The shift and the complete surviving intervals}
Normalize the shift using its nearest multiple of $\ell$:
\begin{equation}\label{eq:signed}
 s=b\ell+\varepsilon v,\qquad 0\leq v\leq t,\qquad
 \varepsilon\in\{1,-1\},\qquad b\geq0.
\end{equation}
If $s=b_0\ell+u$ and $u\leq t$, take $(b,\varepsilon,v)=(b_0,1,u)$.
If $u>t$, take $(b,\varepsilon,v)=(b_0+1,-1,\ell-u)$.
The constraint $s\leq h=A\ell+t$ implies $b\leq A$.

A term of $L_{m,s}$ with $k=a\ell+r$ can be nonzero at $\zeta$ only if
\[
 r\leq t,\qquad \langle r-\varepsilon v\rangle_\ell\leq t,
 \qquad \langle r+\varepsilon v\rangle_\ell\leq t,
\]
where $\langle\cdot\rangle_\ell$ denotes least nonnegative residue.
These conditions are exactly
\begin{equation}\label{eq:interval}
 v\leq r\leq t-v.
\end{equation}
Indeed, if $r<v$, then $\langle r-v\rangle_\ell=\ell+r-v>t$.
If $r\geq v$, both $r-v$ and $r+v$ lie between zero and $2t=\ell-1$,
so the upper condition is simply $r+v\leq t$. If $2v>t$, no terms survive.
This covers both signs of $\varepsilon$.

For a surviving term there are no carries in the two shifted indices:
\[
 k-s=(a-b)\ell+(r-\varepsilon v),\qquad
 k+s=(a+b)\ell+(r+\varepsilon v).
\]
The original bounds $s\leq k\leq h$ now give precisely $b\leq a\leq A$,
and each such block contains the entire interval \eqref{eq:interval}.
For the lower endpoint, $a<b$ is impossible since $r\leq t$ and
$s=b\ell\pm v$; for $a=b$, the condition $r\geq v$ suffices for either
sign. The upper endpoint is $h=A\ell+t$, which does not truncate
\eqref{eq:interval}. Thus there is no discarded boundary block.

Applying Lemma~\ref{lem:blocks} therefore gives the exact root identity
\begin{align}\label{eq:Lreduce}
 L_{m,s}(\zeta)
 &=\left(\sum_{a=b}^{A}\beta_{a-b}\beta_{a+b}\beta_a^2\right)
 \sum_{r=v}^{t-v}[4r+1]_\zeta
 R_{r-v}(\zeta)R_{r+v}(\zeta)R_r(\zeta)^2P_r(\zeta)\notag\\
 &=\left(\sum_{a=b}^{A}\beta_{a-b}\beta_{a+b}\beta_a^2\right)
 L_{\ell,v}(\zeta)=0.
\end{align}
The second equality holds because the omitted terms $t-v<r\leq t$
in $L_{\ell,v}$ contain $R_{r+v}(\zeta)=0$; the final equality is
\eqref{eq:smallzero}. If $2v>t$, both sums in question vanish termwise.
Every term of the original large sum is regular by \eqref{eq:valuations}.
Consequently $L_{m,s}$, as well as $V_{m,s}$, is divisible by every
$\Phi_\ell(q)$ with $\ell\mid m$ and $\ell>1$.

\section{Combining multiplicities and scope of specialization}
Over $K[q]$, the cyclotomic factors are pairwise coprime and
\[
 [m]=\prod_{\substack{\ell\mid m\\\ell>1}}\Phi_\ell(q).
\]
Sections 2--4 show that $L_{m,s}-V_{m,s}$ has order at least one at each
factor of $[m]$. The established Theorem 1 of \cite{KJ} separately supplies
order at least three at $\Phi_m$. Combining these orders yields
\[
 \Phi_m(q)^3\prod_{\substack{\ell\mid m\\1<\ell<m}}\Phi_\ell(q)
 =[m]\Phi_m(q)^2.
\]
The denominator is a unit at every such factor after reduction, by the
regularity proofs above. This completes the proposed proof of
Theorem~\ref{thm:main}.\hfill$\square$

The cases $s=0$ and $s=h$ are included. At $s=h$ both original sides
are actually identical as rational functions, since each sum has one
term and $m+2s=4s+1$. No limit in $m$, primality restriction, or analytic
convergence assertion is used.

The statement proved by the argument is the formal two-parameter
congruence. It immediately admits parameter specializations preserving
a denominator coprime to the modulus. The proof of periodicity for $P_k$
uses generic parameters. At, for example, $c=d=1$, factors
$1-c^\ell$ and $1-c^{-\ell}$ both vanish, so one must not simply substitute
into that intermediate calculation and assert the same block constants.
Although finite exact tests below include $c=d=1$, we do not claim they
prove every such singular specialization for all $m$. If the intended
interpretation of the original question demands all specialized rational
functions, including newly singular parameter values, that stronger
interpretation remains outside this submission's proved scope.

\section{Exact computations and audit}
The script \texttt{exact\_check.py} works over exact rational polynomial
rings. For each cyclotomic factor it first removes its multiplicity from
individual numerator and denominator factors, and only then inverts the
remaining unit modulo $\Phi_\ell^e$. Thus it does not invert a vanishing
$q$-factorial or treat a $0/0$ expression as a numerical value.
The resulting \texttt{exact\_checks.json} records each tested $(m,s,c,d,\ell,e)$,
the zero congruence residual, and minimum term valuations.
In total, 224 exact local checks passed. They include all shifts for
$m=3,5,7,9,11,15$ at $\Phi_m^3$, and every proper divisor greater than
one for $m=9,15,21,25,27$, at both $(c,d)=(2,3)$ and $(1,1)$. Both
sides separately vanish modulo each tested proper-divisor factor.
These finite tests are not part of the proof.

The source and its later related papers were examined for overlap.
Karmakar--Jana's later arXiv:2610.15050v1~\cite{KJlater} proves different
$(\mathrm{G}.2)$ conjectures of Tang and Wang--Xu. Its Section 2 already
contains the same methodological architecture: Pochhammer block limits,
shift-residue restrictions, reflection of surviving intervals, and
combination of cyclotomic factors. The present argument adapts that
strategy to base $q^2$ and the extra parameter factor $P_k$, using the
known smaller instance in place of a separate reflection formula.
This is \emph{not a new general method}; G5 is not passed.
Karmakar's arXiv:2610.33249~\cite{Klater} concerns different base-three
sums and does not state the exact selected conjecture. Neither bounded
review nor keyword search proves priority or continued open status.

The supplied corpus consists of 190 TEX files at commit
\texttt{5887cdea4898f0a3aa5b}\allowbreak\texttt{f7a8e0829697bd2b0f4b}.
Searches found no hits for the selected identifier, Karmakar, or the later
paper's identifier. Generic root-of-unity and block terminology is common
in the corpus, so absence of exact names is not a semantic novelty proof.
Details and retrieval records are in \texttt{audit.txt} and
\texttt{corpus\_audit.json}.

\subsection*{Version-control exclusion: Ni's former Conjecture 5.1}
The first candidate considered for this round was Ni's Conjecture 5.1 in
arXiv:2607.04583v1. Reading the actual version 2~\cite{Ni} showed that the
same assertion is already Theorem 1.4, equation (1.14), with its proof
completed on page 21. It was excluded, not rediscovered. In contrast,
Ni's distinct Conjecture 3.1 remains explicitly a conjecture on page 13
of v2. That is the prior round's target and is not used in this argument.

\begin{center}\small
\begin{tabular}{@{}p{.08\linewidth}p{.84\linewidth}@{}}\toprule
Gate&Assessment\\\midrule
G1&Existing named conjecture: Karmakar--Jana, Conjecture 1.\\
G2&Original equation (1.14), all $m,s$ and formal parameters retained; original pages checked.\\
G3&Conjectural in latest located v1; bounded later review found no exact resolution. Not exhaustively certified.\\
G4&Full formal-parameter proof candidate; independent review pending. No blanket claim for singular specializations.\\
G5&Not passed. The central block method already occurs in the authors' later paper on a different conjecture.\\
G6&Relevant primary follow-up read and method overlap disclosed; corpus and public searches are bounded.\\\bottomrule
\end{tabular}
\end{center}

\begin{thebibliography}{9}
\bibitem{KJ} L. Karmakar and A. Jana, \emph{$q$-Supercongruences Modulo the
Third and Fifth Powers of a Cyclotomic Polynomial}, arXiv:2610.11778v1,
8 October 2026. Theorem 1, equation (1.14), p.~3; Conjecture 1, p.~11.
\url{https://arxiv.org/abs/2610.11778v1}.
\bibitem{KJlater} L. Karmakar and A. Jana, \emph{Proof of Two Conjectural
$q$-Supercongruences Related to Van Hamme's $(\mathrm{G}.2)$ Supercongruence},
arXiv:2610.15050v1, 10 October 2026, especially Section 2.
\url{https://arxiv.org/abs/2610.15050v1}.
\bibitem{Klater} L. Karmakar, \emph{Two parametric $q$-supercongruences
from Andrews' multi-series transformation}, arXiv:2610.33249v1,
23 October 2026. \url{https://arxiv.org/abs/2610.33249v1}.
\bibitem{Ni} H.-X. Ni, \emph{Supercongruences Arising from Truncated
Appell Series}, arXiv:2607.04583v2, 10 July 2026.
\url{https://arxiv.org/abs/2607.04583v2}.
\end{thebibliography}
\end{document}
